% =============================================================================
%  SECTION 6 — Synthesis & Cross-Cutting Insights
%  Slides 6.1–6.3  |  Target time: ~5 minutes
%
%  Corresponds to dissertation Ch. 6 §§2–3
%
%  Drop this \section{} call and the three frames that follow into main.tex
%  after the §5 C3 slides (slide 5.6), before §7 (Conclusions).
% =============================================================================

%%==========================================================================
%% SYNTHESIS
%%==========================================================================
% -----------------------------------------------------------------------------
% Slide 6.2 — Four Cross-Cutting Insights
% -----------------------------------------------------------------------------
\begin{frame}{Through-lines}

    \vspace{-0.3em}
    \begin{enumerate}
        \item \textbf{Representation v. Architecture} \\
            \small Simpler representational solutions may outperform complex
            architectures when the invariance mechanism matches the
            problem structure.
            \hfill\emph{(C2, C3)}
        \note[item]{Rep over architecture: moreover }

        \vspace{0.4em}
        \item \textbf{Task Structure as Critical Moderator} \\
            \small Coordination demands, reward
            sparsity, and heterogeneity type determine what works.
            \hfill\emph{(C1)}

        \vspace{0.4em}
        \item \textbf{Unified Sharing Without Interchangeability} \\
            \small Semantic decomposability unlocks full parameter sharing
            across structurally heterogeneous agents without assuming
            agents are identical.
            \hfill\emph{(C2, C3)}

        \vspace{0.4em}
        \item \textbf{Emergent Robustness} \\
            \small Deployment flexibility — sensor dropout, team changes,
            novel compositions — comes for free from well-designed
            representations, not from explicit robustness objectives.
            \hfill\emph{(C2, C3)}
    \end{enumerate}

    \note{About 1.5 minutes.
    Walk through each insight briskly — one sentence of elaboration each.
    ``First: representation beats architecture. We saw this directly in C3,
    and C2 laid the groundwork.
    Second: task structure matters. C1 showed that curriculum benefits
    depend on coordination demands and heterogeneity type —
    there is no one-size-fits-all.
    Third: parameter sharing does not require agent interchangeability.
    Semantic decomposability is the key condition — if observation
    elements mean the same thing across agents, you can share everything.
    Fourth: robustness emerged for free. We never trained for sensor dropout
    or novel team compositions — the representation handled it naturally.''
    These four points are the take-home value for anyone building
    real heterogeneous systems. Pause briefly before moving on.}
\end{frame}


% -----------------------------------------------------------------------------
% Slide 6.3 — Practitioner Decision Framework
% -----------------------------------------------------------------------------
\begin{frame}[shrink=5]{Practitioner Decision Framework}

    \vspace{2em}
    \centering
    \renewcommand{\arraystretch}{1.5}
    \small
    \begin{tabular}{@{} p{4.5cm} p{5.5cm} @{}}
        \toprule
        \textbf{If your system has \ldots}
            & \textbf{Then consider \ldots} \\
        \midrule
        Behavioral heterogeneity \newline
        {\scriptsize (identical structure, divergent roles)}
            & Adding curriculum-based upsampling into your
            MARL/HARL \hfill\emph{(C1)} \\
        Intrinsic heterogeneity \newline
        {\scriptsize (differing sensor suites)}
            & Implicit Indication \hfill\emph{(C2)} \\
        Rich relational agent interactions \newline
        {\scriptsize (agent-to-agent dependencies)}
            & Architectural augmentation (GNN) — but verify
              the graph carries meaningful structure
              \hfill\emph{(C3)} \\
        Semantically decomposable observations
            & Implicit Indication as the \textbf{default}; add
              architectural complexity only if needed
              \hfill\emph{(C2, C3)} \\
        \bottomrule
    \end{tabular}

    % \vspace{0.5em}
    % \begin{alertblock}{Central Thesis — Revisited}
    %     Training efficiency, heterogeneity handling, and deployment flexibility
    %     in \gls{harl} are fundamentally \textbf{representational} problems
    %     before they are architectural ones.
    % \end{alertblock}

    \note{About 1.5 minutes.
    Frame this as practical advice:
    ``If you are building a heterogeneous multi-agent system,
    here is the decision framework this dissertation provides.''
    Walk down the table row by row.
    ``Behavioral heterogeneity — agents look the same but act differently —
    start with curriculum pretraining.
    Structural heterogeneity — agents have different sensors —
    build a homogenized observation space and let implicit indication
    handle the rest.
    If your agents interact richly with each other,
    graph-based methods may add value — but C3 shows you should verify
    that the graph structure actually carries information.
    And if your observations are semantically decomposable,
    implicit indication should be your default starting point.''
    Then read the alertblock — this is the thesis statement from slide 1.5,
    now supported by all three contributions.
    Let it land. Pause. Then transition:
    ``With those insights established, let me close with
    a summary of what was accomplished and where the work goes next.''}
\end{frame}

